\section{Part III --- The calibration family}\label{sec:calibration}

\subsection{Calibrations}\label{sec:calibration:intro}

A \emph{calibration} records a detector class, the kind of obstruction
it reads, together with the kind of carrier on which that detector is
meaningful.  A low-degree polynomial detector, a mutual-information
detector, a stable local algorithm and a Nullstellensatz certificate
have different native carriers and different notions of descent, and a
transport between two of them is a theorem, not a change of notation.
Part~III records these distinctions in a small typed vocabulary:
calibration records, a criterion for when a new record merely relabels
an old one, records of bridges between calibrations, a table of covered
regions, a four-axis comparison of types, and records of open
transports called \emph{functorial gates}.

The numbered results of this section are elementary facts about
finite records of labels: each one checks a property of declared finite
data, and the names they carry describe the intended use of the data
rather than the depth of the proof.  We state them because they make the catalogue precise and
because the functorial gates of \Cref{sec:calibration:functorial-gates}
are the open problems the catalogue points to.  They prove nothing
about how well a detector performs, about the existence of a bridge, or
about inclusions between complexity classes; region names such as
\(\mathsf{NC}^{1}\) are labels.

\subsection{The seed catalogue}\label{sec:calibration:structure}

\begin{definition}[Calibration]\label{def:calibration}
A \emph{calibration} is a record
\[
  \mathsf{Calibration}
  =
  \{\mathsf{name},\ \mathsf{detectorClass},\ \mathsf{nativeCarrierType}\}
\]
of three labels: its name, the class of detector it uses, and the type
of carrier on which the detector is meaningful.
\end{definition}

The seed catalogue \(\mathcal C_{0}\) is the list of thirteen records in
\Cref{tab:calibration:seed}; ``seed'' means that it is a starting list
to which further calibrations can be added.  It is introduced in this
paper; its
entries name detector classes that have been considered in connection
with the P-versus-NP application, and they are not results of that
application.  Background for the individual detector classes includes
\cite{Lovasz1978Kneser} for the topological line,
\cite{MulmuleySohoni2001GCT} for geometric complexity theory,
\cite{Krajicek2019ProofComplexity} for proof complexity and
Nullstellensatz certificates, and
\cite{AlonKrivelevichSudakov1998PlantedClique,Kearns1998SQ} for planted
problems and statistical queries.

\begin{theorem}[Seed catalogue size]\label{thm:calibration-family-structure}
The seed catalogue \(\mathcal C_{0}\) is a list of thirteen calibration
records; in particular \(|\mathcal C_{0}|\ge 13\).
\end{theorem}

\begin{proof}
The list in \Cref{tab:calibration:seed} has thirteen entries, each with the three fields of \Cref{def:calibration}; the last entry is a
transport between calibrations rather than a detector of its own.
\end{proof}

The catalogue is a seed, not a closed list: the record type is open,
and further calibrations can be adjoined, subject to the rebadge
criterion below.  We call the family \emph{non-descending} in the
intended sense that no entry is presented as a quotient or
specialization of another; this is a description of the entries, not a
proved property.

\begin{table}[H]
  \centering
  \caption{The seed catalogue of \Cref{thm:calibration-family-structure}.}
  \label{tab:calibration:seed}
  \scriptsize
  \renewcommand{\arraystretch}{1.06}
  \begin{tabularx}{\textwidth}{@{}>{\raggedright\arraybackslash}p{0.10\textwidth}>{\raggedright\arraybackslash}p{0.26\textwidth}>{\raggedright\arraybackslash}X>{\raggedright\arraybackslash}p{0.26\textwidth}@{}}
    \toprule
    Symbol & Name & Native carrier type & Detector class \\
    \midrule
    \(\Cvobs\) & volumetric observables & \(AC^0\) / \(AC^0[p]\)-adjacent carriers & low-degree polynomial detector \\
    \(\CMI\) & information calibration & bounded-information packages & mutual information / KL divergence \\
    \(\CRE\) & stability calibration & random-restriction / coupling-stable carriers & stable, local, overlap algorithms \\
    \(\CCh\) & cohomological calibration & cohomology-lifted carriers with Lov\'asz--Kneser reach & topological obstruction \\
    \(\CQamp\) & amplitude calibration & black-box and query carriers & quantum amplitude / query detector \\
    \(\CIalg\) & algebraic calibration & proof-system and certificate-degree carriers & ideal / Nullstellensatz detector \\
    \(\CRreal\) & realizability calibration & lambda-calculus / typed-program carriers & realizability detector \\
    \(\COCstar\) & \(C^*\)-operator calibration & operator-algebraic carriers & \(C^*\)-operator detector \\
    \(\CFexistsSO\) & existential-second-order calibration & logical / locality carriers & descriptive-complexity detector \\
    \(\CSHO\) & higher-order strategy calibration & higher-order game carriers & game-semantic strategy detector \\
    \(\CPPH\) & persistent-homology calibration & scale-dependent topological carriers & persistent-homology detector \\
    \(\CAVNP\) & arithmetic-complexity calibration & GCT / polynomial-family carriers & arithmetic-complexity / VNP detector \\
    \(\CMmot\) & computational motive (META) & cross-calibration transport carriers & motive-level transport detector \\
    \bottomrule
  \end{tabularx}
\end{table}


\subsection{Rebadges}\label{sec:calibration:rebadge}

A new name does not make a new calibration.  A candidate must differ
from every existing entry in the structural data by which a calibration
earns its place.

\begin{definition}[Calibration four-axis profile]\label{def:four-axis-profile}
The \emph{\fourAxis} of a calibration is the quadruple of labels
\[
  p=(p_{\mathsf{alg}},p_{\mathsf{det}},p_{\mathsf{desc}},p_{\mathsf{orient}})
\]
recording its algebraic identities, detector class, descent invariants
and carrier orientation.  A profile \(p\) is a \emph{rebadge} of \(q\),
\(\mathsf{Rebadge}(p,q)\), when \(p=q\) componentwise, and \(p\) is
\emph{genuinely new} relative to a family \(\mathcal F\),
\(\mathsf{GenuineNew}(p,\mathcal F)\), when it is a rebadge of no member
of \(\mathcal F\).
\end{definition}

\begin{theorem}[Rebadge Discipline]\label{thm:rebadge-discipline}
A profile \(p\) is genuinely new relative to \(\mathcal F\) if and only
if, for every \(q\in\mathcal F\), \(p\) and \(q\) differ on at least one
of the four axes.
\end{theorem}

\begin{proof}
\(\mathsf{Rebadge}(p,q)\) is the conjunction of the four componentwise
equalities, so its negation is the disjunction of the four
inequalities; quantify over \(q\in\mathcal F\).
\end{proof}

The criterion compares labels.  Whether two detectors with different
labels are mathematically equivalent is a separate question, which the
criterion does not decide.

\subsection{Bridges between calibrations}\label{sec:calibration:bridges-as-theorems}

\begin{definition}[Calibration-side typed bridge]\label{def:calibration-side-typed-bridge}
A \emph{\TypedBridgeCalibration} is a record
\(B=\{\mathsf{source},\mathsf{target},\mathsf{kind}\}\) consisting of a
source calibration, a target calibration and a kind.\footnote{This is
distinct from the carrier-side bridge of
\Cref{def:carrier-side-typed-bridge}.}
\end{definition}

\begin{definition}[Calibration-bridge kind dichotomy]\label{def:bridge-kind}
The kind of a calibration bridge is one of two values,
\(\mathsf{BridgeKind}::=\mathsf{basisChange}\mid\mathsf{substantive}\):
a \emph{basis change} relates two presentations of the same detector and
carrier content, and a \emph{substantive} bridge asserts a theorem
relating different invariants.
\end{definition}

\begin{theorem}[Bridge kinds]\label{thm:bridges-as-theorems}
Every calibration bridge record has kind \(\mathsf{basisChange}\) or
kind \(\mathsf{substantive}\).
\end{theorem}

\begin{proof}
\(\mathsf{BridgeKind}\) has exactly these two values.
\end{proof}

The content of the theorem is a convention: a transport between
calibrations is recorded only together with a declared kind, and is
then either a change of basis or a theorem that has to be proved.
Placing two calibrations side by side without such a declaration does
not produce a bridge record.  The kind label itself does not supply
the change of basis or the proof.

\subsection{Direction}\label{sec:calibration:bridge-type-direction}

Transports between calibrations usually have a direction: worst case to
average case, nonuniform to uniform, statistical to algorithmic.

\begin{definition}[Calibration-bridge reversal]\label{def:calibration-bridge-reversal}
The \emph{reversal} of a bridge record \(B=(S,T,k)\) is
\(B^{\mathrm{op}}=(T,S,k)\).
\end{definition}

\begin{theorem}[Bridge Type-Direction]\label{thm:bridge-type-direction}
For every calibration bridge record \(B\),
\(B=B^{\mathrm{op}}\) if and only if
\(\mathsf{source}(B)=\mathsf{target}(B)\).
\end{theorem}

\begin{proof}
If \(B=B^{\mathrm{op}}\), comparing source fields gives \(S=T\);
conversely, if \(S=T\) the two records coincide.
\end{proof}

Thus a bridge record between different calibrations is never its own
reversal.  Whether a bridge \(S\to T\) yields a usable bridge
\(T\to S\) is a mathematical question about the transport, which the
records do not model.

\subsection{Coverage}\label{sec:calibration:coverage-family-geometry}

Each calibration \(c\) is assigned a set \(\operatorname{Cov}(c)\) of
region labels, and the family coverage is
\(\operatorname{Cov}(\mathcal C_{0})=\bigcup_{c\in\mathcal C_{0}}\operatorname{Cov}(c)\).
The labels and the assignment form a finite declared table, shown in
\Cref{tab:coverage-regions}; the assignment is a modelling choice, and
the seven calibrations not listed there are assigned no region.  Two uncovered labels are singled out as
\emph{typed coverage holes}: \(\xiNCone\), with region
\(\mathsf{NC}^{1}\), and \(\xiAdaptive\), with region
\(\mathsf{adaptivePolyTime}\).  The standard references for the classes
suggested by the uncovered labels are \cite{Cook1985NCTaxonomy} for the
\(\mathsf{NC}\) hierarchy and \cite{KarpLipton1982Advice} for
\(\mathsf{P/poly}\).

\begin{table}[H]
  \centering
  \caption{The declared coverage table of \Cref{thm:coverage-family-geometry}.  The assignment of regions to calibrations is a modelling choice; calibrations not listed are assigned no region.}
  \label{tab:coverage-regions}
  \small
  \renewcommand{\arraystretch}{1.08}
  \begin{tabularx}{\textwidth}{@{}>{\raggedright\arraybackslash}p{0.30\textwidth}>{\raggedright\arraybackslash}X>{\raggedright\arraybackslash}p{0.22\textwidth}@{}}
    \toprule
    Region label & Assigned to & Status \\
    \midrule
    bounded-depth & \(\Cvobs\), \(\CMI\) & covered \\
    low-degree & \(\Cvobs\) & covered \\
    stable-local & \(\CRE\) & covered \\
    query & \(\CQamp\) & covered \\
    topological & \(\CCh\) & covered \\
    low-Nullstellensatz-degree & \(\CIalg\) & covered \\
    \(\mathsf{NC}^{1}\) & none & uncovered; hole \(\xiNCone\) \\
    \(\mathsf{NC}^{2}\) & none & uncovered \\
    \(\mathsf{P/poly}\) & none & uncovered \\
    adaptive polynomial time & none & uncovered; hole \(\xiAdaptive\) \\
    \bottomrule
  \end{tabularx}
\end{table}


\begin{theorem}[Declared coverage table]\label{thm:coverage-family-geometry}
In the declared table: the labels bounded-depth, low-degree,
stable-local, query, topological and low-Nullstellensatz-degree belong
to \(\operatorname{Cov}(\mathcal C_{0})\); the labels \(\mathsf{NC}^{1}\),
\(\mathsf{NC}^{2}\), \(\mathsf{P/poly}\) and adaptive polynomial time
do not; the regions of \(\xiNCone\) and \(\xiAdaptive\) are uncovered;
two distinct seed calibrations, \(\Cvobs\) and \(\CMI\), both cover the
bounded-depth label; and for every family \(\mathcal F\) and every
calibration \(c\) with \(\operatorname{Cov}(c)=\varnothing\),
\[
  \operatorname{Cov}(\mathcal F\cup\{c\})=\operatorname{Cov}(\mathcal F).
\]
\end{theorem}

\begin{proof}
The membership statements are checked entry by entry in the table.  The
last statement holds because the union with an empty set is unchanged.
\end{proof}

The theorem verifies the table; it proves no inclusion, exclusion or
separation between the complexity classes named by the labels.  Its
last clause is a reminder that adding names to the catalogue does not
by itself enlarge its coverage: a new calibration changes the coverage
only if it covers a new region, or if a bridge theorem carries content
into one.

\subsection{Native type and target type}
\label{sec:calibration:framework-type-self-char}

\begin{definition}[Framework-type]\label{def:framework-type}
A \emph{\FType} is a record of four labels: uniformity, worst-caseness,
scoping and packaging.  Two framework-types \emph{match} when they agree
on all four.  For the P-versus-NP application we record
\[
\begin{aligned}
  F_{\mathrm{PvNP}}
  &=(\mathsf{nonuniform},\mathsf{averageCase},\mathsf{scoped},\mathsf{packageBased}),\\
  T_{\mathrm{Clay}}
  &=(\mathsf{uniform},\mathsf{worstCase},\mathsf{unconditional},\mathsf{languageBased}).
\end{aligned}
\]
\end{definition}

The first two axes correspond to the distinctions surveyed in
\cite{Impagliazzo1995FiveWorlds}.  The tuple \(F_{\mathrm{PvNP}}\)
describes the calibration catalogue of this section, not every
construction of the application: the P-versus-NP paper itself proves a
uniform, worst-case machine translation and a conditional separation of
standard languages.  The comparison below concerns the two recorded
tuples and does not classify that proof.

\begin{theorem}[Recorded type mismatch]\label{thm:framework-type-self-char}
The recorded types \(F_{\mathrm{PvNP}}\) and \(T_{\mathrm{Clay}}\) do not
match, and the gate catalogue of \Cref{cat:named-functorial-gates} has
three entries.
\end{theorem}

\begin{proof}
The first labels, \(\mathsf{nonuniform}\) and \(\mathsf{uniform}\), are
different (so are the other three).  The count is
\Cref{cat:named-functorial-gates}.
\end{proof}

\begin{schema}[Framework-general candidate]\label{sch:framework-type-pending}
For every formed closure \(c\) in the trace-state-only column, the
native framework-type assigned to \(c\) and the framework-type of its
standard target do not match.
\end{schema}

Both type assignments in the schema are primitive.  The schema is not
used in the proof of \Cref{thm:framework-type-self-char}; it records the
conjectured extension beyond the P-versus-NP case.  The mismatch says
where the catalogue's native reach lies, and therefore which transports
would have to be supplied to reach the standard target; the next
subsection names three of them.

\subsection{Functorial gates}\label{sec:calibration:functorial-gates}

\begin{definition}[Named functorial gate]\label{def:functorial-gate}
A \emph{named functorial gate} is a record
\[
  \{\mathsf{name},\mathsf{source},\mathsf{target},\mathsf{direction},\mathsf{theoremObligation}\}
\]
of labels: a typed source, a typed target, a direction, and the
statement of the theorem that a transport from source to target would
have to prove.  The record does not assert that the transport exists.
\end{definition}

\begin{table}[H]
  \centering
  \caption{The three named functorial gates of \Cref{cat:named-functorial-gates}.  Each records a transport that has not been proved.}
  \label{tab:functorial-gates}
  \small
  \renewcommand{\arraystretch}{1.1}
  \begin{tabularx}{\textwidth}{@{}>{\raggedright\arraybackslash}p{0.08\textwidth}>{\raggedright\arraybackslash}p{0.22\textwidth}>{\raggedright\arraybackslash}p{0.22\textwidth}>{\raggedright\arraybackslash}X@{}}
    \toprule
    Gate & Source & Target & Obligation \\
    \midrule
    \(\GateFunctionDegree\) & degree of a low-degree detector (\(\Cvobs\)) & degree of a Nullstellensatz certificate & relate the degree seen by the detector to the degree of an ideal-membership certificate \\
    \(\GatePlantedStatistical\) & average-case (planted) statistical detection & worst-case hardness of an NP-hard language & carry detection content from the average-case side to the worst case; as a hardness statement, show that the language being hard in the worst case makes the detection problem hard on average \\
    \(\GateUniformTime\) & nonuniform families & uniform time hierarchy & realize a nonuniform family in a uniform time-hierarchy argument of the kind used by Williams \cite{Williams2014NonUniformACC} \\
    \bottomrule
  \end{tabularx}
\end{table}


\begin{theorem}[Gate catalogue]\label{cat:named-functorial-gates}
The gate catalogue \(\mathcal G_{\mathrm{named}}=[G_{1},G_{2},G_{3}]\)
of \Cref{tab:functorial-gates} has three entries, and each has a
nonempty name and a nonempty theorem obligation.
\end{theorem}

\begin{proof}
Inspection of the three records.
\end{proof}

None of the three obligations is proved here; they are the open
problems of \Cref{sec:predictions:functorial-gates}.  The gate records
complement the classical barrier results---relativization, natural
proofs and algebrization
\cite{BakerGillSolovay1975Relativization,RazborovRudich1997NaturalProofs,AaronsonWigderson2009Algebrization}---which
show that certain families of techniques cannot cross certain
boundaries; a gate states what a crossing would have to prove.  The
P-versus-NP application also contains conditional results restricting
feature-based proof attempts, under explicit hypotheses: a
feature-decomposition record, a joint-access operation, and assumptions
that expressions of given classes do not prove the target
\cite{Tsiokos2026PvNP}.  These results do not constitute a barrier of
the scope of the classical ones.  The three gates are examples, not a complete list.
